\chapter{Unix/Linux for Scientific Computing}

Modern astronomy relies heavily on large datasets, remote servers,
automated pipelines, and high-performance computing systems. Most
scientific computing infrastructure in astronomy runs on Unix/Linux.
For this reason, learning to use the Linux command line is an essential
skill for modern astronomers.

This chapter introduces the most important Unix/Linux commands needed
for scientific work. Instead of treating Linux as a list of isolated
commands, we will emphasize practical usage in a research environment:
navigating directories, managing files, inspecting data, searching for
information, connecting to remote servers, and running long computations.

\section{Why Linux Matters in Astronomy}

Astronomy is a data-intensive science. A single observing run may
produce thousands of FITS files, while a numerical simulation may create
many gigabytes of output. Graphical interfaces are often convenient for
small tasks, but they become inefficient when the number of files grows
large.

Linux is especially useful in astronomy because it allows researchers to:

\begin{itemize}
\item manage large collections of files efficiently,
\item automate repetitive tasks,
\item process data on remote servers,
\item combine many small tools into a powerful workflow,
\item interact with shared computing clusters.
\end{itemize}

\begin{defbox}[Scientific Computing Environment]
In astronomy, Linux is not merely another operating system. It is the
standard working environment for data reduction, simulations, remote
computing, and pipeline development.
\end{defbox}

\begin{examplebox}[A Typical Research Task]
Suppose a telescope produces a large number of FITS images in a single
night. A researcher may need to:

\begin{itemize}
\item log into a remote server,
\item create directories for raw and processed data,
\item move FITS files into the correct folder,
\item inspect the first few lines of a log file,
\item search for a target observation,
\item compress and archive the final results.
\end{itemize}

Linux provides simple tools for each of these steps.
\end{examplebox}

\section{The Command Line Interface}

A graphical user interface (GUI) allows a user to click icons and open
menus. A command-line interface (CLI) allows the user to type commands
directly in a terminal.

\begin{defbox}[Command-Line Interface]
A command-line interface is a text-based way of interacting with the
operating system. Instead of selecting actions from menus, the user
enters commands explicitly.
\end{defbox}

The CLI is especially valuable in scientific computing because it is:

\begin{itemize}
\item efficient,
\item scriptable,
\item flexible,
\item well suited for remote access.
\end{itemize}

For example, moving many FITS files into a directory is a one-line task:

\begin{lstlisting}
$ mkdir raw_data
$ mv *.fits raw_data/
\end{lstlisting}

\section{The Linux File System}

Linux organizes files using a hierarchical tree structure.

\begin{defbox}[Root Directory]
The top of the Linux file system is called the root directory, written
as \texttt{/}. Every file and directory belongs somewhere under this root.
\end{defbox}

For example:

\begin{lstlisting}
/home/student/project/data
\end{lstlisting}

Here:

\begin{itemize}
\item \texttt{/} is the root directory,
\item \texttt{home} contains user home directories,
\item \texttt{student} is the user's home directory,
\item \texttt{project/data} are subdirectories.
\end{itemize}

Some common directories are:

\begin{itemize}
\item \texttt{/home}: user home directories
\item \texttt{/bin}: basic executable programs
\item \texttt{/etc}: configuration files
\item \texttt{/tmp}: temporary files
\item \texttt{/usr}: installed software and libraries
\end{itemize}

\subsection{Absolute and Relative Paths}

An \textbf{absolute path} starts from the root directory:

\begin{lstlisting}
/home/student/project/data
\end{lstlisting}

A \textbf{relative path} starts from the current working directory:

\begin{lstlisting}
project/data
\end{lstlisting}

Understanding the difference between absolute and relative paths is one
of the most important early steps in learning Linux.

\section{Navigating Directories}

\subsection{\texttt{pwd}: Print Working Directory}

The command \texttt{pwd} shows the full path of the current directory.

\begin{lstlisting}
$ pwd
\end{lstlisting}

Example output:

\begin{lstlisting}
/home/student/project
\end{lstlisting}

This command is especially useful when you are working on a remote server
and want to confirm your location before copying or deleting files.

\subsection{\texttt{ls}: List Directory Contents}

The command \texttt{ls} lists files and directories.

\begin{lstlisting}
$ ls
\end{lstlisting}

Some commonly used options are:

\begin{itemize}
\item \texttt{-l}: long listing format
\item \texttt{-h}: human-readable file sizes
\item \texttt{-a}: include hidden files
\item \texttt{-t}: sort by modification time
\item \texttt{-r}: reverse the sorting order
\end{itemize}

Examples:

\begin{lstlisting}
$ ls -l
$ ls -lh
$ ls -la
$ ls -ltr
\end{lstlisting}

The long listing format provides much more information, for example:

\begin{lstlisting}
-rw-r--r-- 1 user group 2.1G catalog.fits
drwxr-xr-x 2 user group 4.0K scripts
\end{lstlisting}

This includes file permissions, number of links, owner, group, file size,
modification time, and file name.

Files beginning with a dot, such as \texttt{.bashrc}, are hidden files.
These are usually configuration files.

\begin{protip}[A Very Common Combination]
The command
\begin{lstlisting}
$ ls -lh
\end{lstlisting}
is one of the most useful ways to inspect files, because it combines a
detailed view with readable file sizes such as \texttt{K}, \texttt{M},
and \texttt{G}.
\end{protip}

\subsection{\texttt{cd}: Change Directory}

The command \texttt{cd} changes the current working directory.

\begin{lstlisting}
$ cd data
$ cd ..
$ cd ~
$ cd /
\end{lstlisting}

Common patterns:

\begin{itemize}
\item \texttt{cd data}: enter the directory \texttt{data}
\item \texttt{cd ..}: move to the parent directory
\item \texttt{cd ~}: return to the home directory
\item \texttt{cd /}: go to the root directory
\item \texttt{cd -}: return to the previous directory
\end{itemize}

The command \texttt{cd -} is especially helpful when moving back and forth
between two locations.

\begin{protip}[Tab Completion]
When typing a file or directory name, press the \texttt{Tab} key to
auto-complete it. This reduces typing mistakes and speeds up navigation.
\end{protip}

\subsection{\texttt{tree}: Display Directory Structure}

If available on your system, \texttt{tree} shows the directory hierarchy
in a visually clear way.

\begin{lstlisting}
$ tree
\end{lstlisting}

This is helpful when you want to understand the structure of a project.

\section{Managing Files and Directories}

\subsection{\texttt{mkdir}: Make Directories}

The command \texttt{mkdir} creates directories.

\begin{lstlisting}
$ mkdir project
\end{lstlisting}

A very useful option is:

\begin{itemize}
\item \texttt{-p}: create parent directories if needed
\end{itemize}

Example:

\begin{lstlisting}
$ mkdir -p project/data/raw
\end{lstlisting}

This creates multiple nested directories at once.

\subsection{\texttt{cp}: Copy Files and Directories}

The command \texttt{cp} copies files or directories.

Copy one file:

\begin{lstlisting}
$ cp script.py backup_script.py
\end{lstlisting}

Useful options include:

\begin{itemize}
\item \texttt{-r}: copy directories recursively
\item \texttt{-i}: ask before overwriting
\item \texttt{-v}: show what is being copied
\end{itemize}

Examples:

\begin{lstlisting}
$ cp -i report.txt backup.txt
$ cp -r old_run new_run
$ cp -rv source_dir target_dir
\end{lstlisting}

For beginners, \texttt{-i} is a very useful safety option because it
helps prevent accidental overwriting.

\subsection{\texttt{mv}: Move or Rename Files}

The command \texttt{mv} can move a file to another directory or rename it.

Rename a file:

\begin{lstlisting}
$ mv old_name.txt new_name.txt
\end{lstlisting}

Move files into a directory:

\begin{lstlisting}
$ mv *.fits raw_data/
\end{lstlisting}

Useful options include:

\begin{itemize}
\item \texttt{-i}: ask before overwriting
\item \texttt{-v}: show what is being moved
\end{itemize}

Examples:

\begin{lstlisting}
$ mv -i draft.txt final.txt
$ mv -v *.log logs/
\end{lstlisting}

\subsection{\texttt{rm}: Remove Files and Directories}

The command \texttt{rm} removes files.

\begin{lstlisting}
$ rm file.txt
\end{lstlisting}

Useful options include:

\begin{itemize}
\item \texttt{-r}: remove directories recursively
\item \texttt{-f}: force removal without asking
\item \texttt{-i}: ask before removing
\end{itemize}

Examples:

\begin{lstlisting}
$ rm -i file.txt
$ rm -r old_directory
$ rm -ri old_project
\end{lstlisting}

\begin{warnbox}[Dangerous Command]
The command
\begin{lstlisting}
$ rm -rf *
\end{lstlisting}
can permanently delete all files in the current directory. In Linux,
deleted files usually do not go to a recycle bin. Always check your
location with \texttt{pwd} and inspect the directory with \texttt{ls}
before using \texttt{rm}.
\end{warnbox}

\section{Viewing and Inspecting Files}

In scientific work, it is often useful to inspect a file before using it
in analysis.

\subsection{\texttt{cat}: Display a File}

The command \texttt{cat} prints the contents of a file.

\begin{lstlisting}
$ cat file.txt
\end{lstlisting}

This is useful for short text files, configuration files, or scripts.
For very large files, \texttt{less} is usually better.

\subsection{\texttt{less}: View a File Interactively}

The command \texttt{less} opens a file in an interactive viewer.

\begin{lstlisting}
$ less large_file.txt
\end{lstlisting}

Inside \texttt{less}, you can scroll, search, and quit with \texttt{q}.

A useful pattern is:

\begin{lstlisting}
$ less run.log
\end{lstlisting}

which lets you inspect a long log file without printing everything to
the terminal at once.

\subsection{\texttt{head}: Show the Beginning of a File}

The command \texttt{head} shows the first lines of a file.

\begin{lstlisting}
$ head catalog.txt
\end{lstlisting}

Useful option:

\begin{itemize}
\item \texttt{-n}: specify the number of lines
\end{itemize}

Example:

\begin{lstlisting}
$ head -20 catalog.txt
\end{lstlisting}

This is very useful for checking whether a file has a header and for
inspecting the general structure of tabular data.

\subsection{\texttt{tail}: Show the End of a File}

The command \texttt{tail} shows the last lines of a file.

\begin{lstlisting}
$ tail run.log
\end{lstlisting}

Useful options include:

\begin{itemize}
\item \texttt{-n}: specify the number of lines
\item \texttt{-f}: follow a growing file
\end{itemize}

Examples:

\begin{lstlisting}
$ tail -50 run.log
$ tail -f run.log
\end{lstlisting}

The option \texttt{-f} is especially useful when monitoring a log file
while a program is still running.

\subsection{\texttt{wc}: Count Lines, Words, and Characters}

The command \texttt{wc} counts lines, words, and characters.

\begin{lstlisting}
$ wc catalog.txt
\end{lstlisting}

Useful options include:

\begin{itemize}
\item \texttt{-l}: count lines only
\item \texttt{-w}: count words only
\item \texttt{-c}: count bytes only
\end{itemize}

Examples:

\begin{lstlisting}
$ wc -l catalog.txt
$ wc -c image_header.txt
\end{lstlisting}

For data analysis, \texttt{wc -l} is often used to estimate how many rows
a text catalog contains.

\subsection{\texttt{file}: Identify File Type}

The command \texttt{file} tries to identify the type of a file.

\begin{lstlisting}
$ file image.fits
$ file archive.tar.gz
\end{lstlisting}

This is useful when the file extension is missing or unreliable.

\section{Searching for Information}

\subsection{\texttt{grep}: Search Text Patterns}

The command \texttt{grep} searches for lines containing a pattern.

\begin{lstlisting}
$ grep SN catalog.txt
\end{lstlisting}

Useful options include:

\begin{itemize}
\item \texttt{-i}: ignore case
\item \texttt{-n}: show line numbers
\item \texttt{-r}: search recursively in directories
\item \texttt{-v}: show lines that do not match
\end{itemize}

Examples:

\begin{lstlisting}
$ grep -i galaxy notes.txt
$ grep -n ERROR run.log
$ grep -r "target_name" project/
$ grep -v "#" data.txt
\end{lstlisting}

The last example is often useful when removing comment lines from a text file.

\subsection{\texttt{find}: Search for Files}

The command \texttt{find} searches for files and directories in a directory tree.

\begin{lstlisting}
$ find . -name "*.fits"
\end{lstlisting}

Useful patterns include:

\begin{itemize}
\item \texttt{-name}: search by name
\item \texttt{-type f}: search only files
\item \texttt{-type d}: search only directories
\end{itemize}

Examples:

\begin{lstlisting}
$ find . -name "*.fits"
$ find . -type f -name "*.log"
$ find . -type d -name "results"
\end{lstlisting}

This command is very important when working with large project directories.

\section{Wildcards, Pipes, and Redirection}

\subsection{Wildcards}

Wildcards allow one command to match many filenames.

\begin{itemize}
\item \texttt{*}: match any sequence of characters
\item \texttt{?}: match exactly one character
\end{itemize}

Examples:

\begin{lstlisting}
$ ls *.fits
$ ls image_?.fits
\end{lstlisting}

Wildcards are essential when dealing with many similar files.

\subsection{Pipes}

A pipe, written as \texttt{|}, sends the output of one command into another.

\begin{lstlisting}
$ cat catalog.txt | grep SN
\end{lstlisting}

This means: first display the file, then filter lines containing
\texttt{SN}.

Another example:

\begin{lstlisting}
$ ls -lh | less
\end{lstlisting}

Here the output of \texttt{ls -lh} is passed into \texttt{less} so that
it can be viewed page by page.

\subsection{Redirection}

Redirection writes command output into a file.

\begin{itemize}
\item \texttt{>}: overwrite a file
\item \texttt{>>}: append to a file
\end{itemize}

Examples:

\begin{lstlisting}
$ grep SN catalog.txt > sn_list.txt
$ grep galaxy catalog.txt >> result.txt
\end{lstlisting}

This is useful when building simple command-line workflows.

\section{Remote Login and File Transfer}

Many astronomical computations are performed on remote servers rather than
on a local laptop.

\subsection{\texttt{ssh}: Secure Shell}

The command \texttt{ssh} logs into a remote machine.

\begin{lstlisting}
$ ssh user@server
\end{lstlisting}

Example:

\begin{lstlisting}
$ ssh wangf@astrocluster.edu
\end{lstlisting}

Common option:

\begin{itemize}
\item \texttt{-X}: enable X11 forwarding on some systems
\end{itemize}

Example:

\begin{lstlisting}
$ ssh -X user@server
\end{lstlisting}

\subsection{\texttt{scp}: Secure Copy}

The command \texttt{scp} copies files between computers.

Upload a file:

\begin{lstlisting}
$ scp data.fits user@server:/home/user/
\end{lstlisting}

Download a file:

\begin{lstlisting}
$ scp user@server:/home/user/result.txt .
\end{lstlisting}

Useful options include:

\begin{itemize}
\item \texttt{-r}: copy directories recursively
\item \texttt{-P}: specify a custom port
\end{itemize}

Examples:

\begin{lstlisting}
$ scp -r results/ user@server:/home/user/
$ scp -P 2222 data.fits user@server:/home/user/
\end{lstlisting}

\subsection{\texttt{ssh-keygen}: Create SSH Keys}

SSH keys allow password-free login.

\begin{lstlisting}
$ ssh-keygen
\end{lstlisting}

This generates a private key and a public key. The public key can be
copied to a server for authentication.

\begin{protip}[Remote Computing Is Standard]
Many students first encounter Linux through remote login to a department
server or cluster. Learning \texttt{ssh}, \texttt{scp}, and SSH keys
early is extremely valuable.
\end{protip}

\section{Permissions}

Every file and directory in Linux has permissions.

A typical permission string looks like:

\begin{lstlisting}
-rwxr-xr--
\end{lstlisting}

This can be read as:

\begin{itemize}
\item first character: file type
\item next three characters: owner permissions
\item next three characters: group permissions
\item last three characters: permissions for others
\end{itemize}

The symbols mean:

\begin{itemize}
\item \texttt{r}: read
\item \texttt{w}: write
\item \texttt{x}: execute
\end{itemize}

\subsection{\texttt{chmod}: Change Permissions}

The command \texttt{chmod} changes permissions.

Examples:

\begin{lstlisting}
$ chmod +x run.sh
$ chmod o-w file.txt
$ chmod u+r script.py
\end{lstlisting}

Here:

\begin{itemize}
\item \texttt{u} means user (owner)
\item \texttt{g} means group
\item \texttt{o} means others
\item \texttt{a} means all
\end{itemize}

\subsection{\texttt{chown}: Change Ownership}

On systems where you have permission, \texttt{chown} changes file ownership.

\begin{lstlisting}
$ chown user file.txt
\end{lstlisting}

This command is more common in system administration than in ordinary
student work, but it is useful to recognize it.

\section{Disk Usage and Storage}

Astronomical datasets can quickly consume large amounts of storage.

\subsection{\texttt{du}: Disk Usage}

The command \texttt{du} shows directory sizes.

\begin{lstlisting}
$ du -sh *
\end{lstlisting}

Useful options include:

\begin{itemize}
\item \texttt{-s}: summary only
\item \texttt{-h}: human-readable sizes
\end{itemize}

This helps identify which directories occupy the most space.

\subsection{\texttt{df}: Disk Free Space}

The command \texttt{df} shows available disk space.

\begin{lstlisting}
$ df -h
\end{lstlisting}

This is important on shared systems, where storage quotas are often limited.

\section{Monitoring and Managing Processes}

Many scientific programs run for a long time, so it is important to
check and manage active processes.

\subsection{\texttt{top}: Monitor Running Processes}

The command \texttt{top} shows running processes in real time.

\begin{lstlisting}
$ top
\end{lstlisting}

It provides information about CPU usage, memory usage, and active tasks.

\subsection{\texttt{ps}: List Processes}

The command \texttt{ps} lists active processes.

\begin{lstlisting}
$ ps
$ ps -u username
$ ps aux
\end{lstlisting}

The version \texttt{ps aux} is widely used because it provides a detailed
overview of running processes.

\subsection{\texttt{kill}: Terminate a Process}

The command \texttt{kill} stops a process by its process ID.

\begin{lstlisting}
$ kill PID
\end{lstlisting}

If necessary, force termination:

\begin{lstlisting}
$ kill -9 PID
\end{lstlisting}

\begin{warnbox}[Use \texttt{kill -9} Carefully]
The option \texttt{-9} forces a process to stop immediately. This can
sometimes leave incomplete output or corrupted temporary files. Use
ordinary \texttt{kill} first whenever possible.
\end{warnbox}

\section{Terminal Productivity Tools}

\subsection{\texttt{history}}

The command \texttt{history} lists previously used commands.

\begin{lstlisting}
$ history
\end{lstlisting}

This is useful when you want to repeat or modify earlier commands.

\subsection{\texttt{clear}}

The command \texttt{clear} clears the terminal screen.

\begin{lstlisting}
$ clear
\end{lstlisting}

\subsection{Command Recall}

In most terminals, the up-arrow key retrieves earlier commands. This is
one of the simplest but most useful command-line habits.

\section{Long-Running Jobs and Persistent Sessions}

Many astronomical calculations run for hours or days. If you are connected
through \texttt{ssh}, a network interruption may stop your work unless you
use a persistent session manager.

\subsection{\texttt{tmux}}

The command \texttt{tmux} starts a persistent terminal session.

\begin{lstlisting}
$ tmux
\end{lstlisting}

Create a named session:

\begin{lstlisting}
$ tmux new -s reduction
\end{lstlisting}

Useful commands:

\begin{itemize}
\item detach from a session: \texttt{Ctrl-b d}
\item list sessions: \texttt{tmux ls}
\item reattach to a session: \texttt{tmux attach -t reduction}
\end{itemize}

\begin{protip}[Why \texttt{tmux} Matters]
If you start a long-running program through \texttt{ssh} and then lose
your connection, the program may stop. Running it inside \texttt{tmux}
protects you from this problem.
\end{protip}

\section{Archiving and Compression}

Scientific projects often contain many files that need to be archived or shared.

\subsection{\texttt{tar}: Archive Files}

Create a compressed archive:

\begin{lstlisting}
$ tar -czvf project.tar.gz project/
\end{lstlisting}

Extract an archive:

\begin{lstlisting}
$ tar -xzvf project.tar.gz
\end{lstlisting}

Common options include:

\begin{itemize}
\item \texttt{-c}: create archive
\item \texttt{-x}: extract archive
\item \texttt{-z}: use gzip compression
\item \texttt{-v}: verbose output
\item \texttt{-f}: specify archive filename
\end{itemize}

\subsection{\texttt{gzip}: Compress Files}

Compress a file:

\begin{lstlisting}
$ gzip large_file.txt
\end{lstlisting}

Decompress it:

\begin{lstlisting}
$ gunzip large_file.txt.gz
\end{lstlisting}

\section{Basic HPC Cluster Commands}

Many astronomy students eventually use a computing cluster. Different
clusters use different schedulers, but one common system is SLURM.

Some useful SLURM commands are:

\begin{lstlisting}
$ squeue
$ sbatch job.sh
$ scancel JOBID
\end{lstlisting}

Meaning:

\begin{itemize}
\item \texttt{squeue}: list queued and running jobs
\item \texttt{sbatch job.sh}: submit a batch script
\item \texttt{scancel JOBID}: cancel a job
\end{itemize}

At this stage, you do not need to master SLURM in detail, but you should
be aware that modern astronomy often relies on job schedulers.

\section{A Typical Astronomy Workflow}

A typical research workflow may look like this:

\begin{enumerate}
\item Log into a remote cluster using \texttt{ssh}.
\item Create a project directory using \texttt{mkdir}.
\item Transfer data using \texttt{scp}.
\item Inspect files using \texttt{ls}, \texttt{head}, and \texttt{wc}.
\item Search for target files using \texttt{find}.
\item Run a long analysis inside \texttt{tmux}.
\item Monitor processes using \texttt{top} or \texttt{squeue}.
\item Archive the final results using \texttt{tar}.
\end{enumerate}

This workflow is much closer to actual scientific practice than learning
commands in isolation.

\section{Command Summary}

\begin{center}
\begin{tabular}{ll}
\textbf{Command} & \textbf{Purpose} \\
\hline
\texttt{pwd}   & show current directory \\
\texttt{ls}    & list files and directories \\
\texttt{cd}    & change directory \\
\texttt{mkdir} & create directories \\
\texttt{cp}    & copy files and directories \\
\texttt{mv}    & move or rename files \\
\texttt{rm}    & remove files and directories \\
\texttt{cat}   & display a file \\
\texttt{less}  & view a file interactively \\
\texttt{head}  & show the beginning of a file \\
\texttt{tail}  & show the end of a file \\
\texttt{wc}    & count lines, words, and bytes \\
\texttt{grep}  & search text patterns \\
\texttt{find}  & search for files \\
\texttt{ssh}   & log into a remote server \\
\texttt{scp}   & transfer files securely \\
\texttt{chmod} & change permissions \\
\texttt{du}    & inspect disk usage \\
\texttt{df}    & check free disk space \\
\texttt{top}   & monitor running processes \\
\texttt{ps}    & list processes \\
\texttt{kill}  & stop a process \\
\texttt{tmux}  & create persistent terminal sessions \\
\texttt{tar}   & create or extract archives \\
\texttt{gzip}  & compress files \\
\end{tabular}
\end{center}

\section{Exercises}

\begin{exercisebox}[Create a Project Directory]
Create the following directory structure in your home directory:

\begin{lstlisting}
astro_project/
    data/
    scripts/
    results/
\end{lstlisting}

Then verify the structure using \texttt{ls} or \texttt{tree}.
\end{exercisebox}

\begin{exercisebox}[Practice File Operations]
Create a text file, copy it, rename it, and move it into another directory.
Write down the commands you used.
\end{exercisebox}

\begin{exercisebox}[Inspect a Data File]
Take a text file or catalog and use the following commands on it:

\begin{lstlisting}
$ head file.txt
$ tail file.txt
$ wc -l file.txt
$ less file.txt
\end{lstlisting}

Explain what information each command provides.
\end{exercisebox}

\begin{exercisebox}[Search for FITS Files]
Use the \texttt{find} command to search for all FITS files in a directory.
Save the output into a file named \texttt{fits\_list.txt}.
\end{exercisebox}

\begin{exercisebox}[Remote Login and Transfer]
Log into a remote machine using \texttt{ssh}. Then transfer one file using
\texttt{scp}. Record both commands.
\end{exercisebox}

\begin{exercisebox}[Disk Usage]
Use \texttt{du -sh *} to inspect the sizes of directories in a project folder.
Then use \texttt{df -h} to check the available disk space on the system.
\end{exercisebox}

\begin{exercisebox}[Persistent Session]
Start a \texttt{tmux} session, create a file inside it, detach from the
session, and then reattach to confirm that the session remains active.
\end{exercisebox}

\begin{exercisebox}[Archive a Project]
Create a compressed archive of a directory using \texttt{tar}. Then extract
it and verify that the files are present.
\end{exercisebox}